% 图契约：本题的源、边界和推导逐项对应；箭头符号按图注解释。
\documentclass[tikz,border=7pt]{standalone}
\usepackage{fontspec,xeCJK,amsmath,amssymb}
\IfFontExistsTF{Times New Roman}{\setmainfont{Times New Roman}}{\setmainfont{TeX Gyre Termes}}
\IfFontExistsTF{Heiti SC}{\setCJKmainfont[AutoFakeBold=1.5]{Heiti SC}}{\setCJKmainfont[AutoFakeBold=1.5]{FandolHei-Regular}}
\usetikzlibrary{arrows.meta,bending,calc,patterns,positioning}
\definecolor{ink}{HTML}{243745}
\definecolor{blue}{HTML}{275DAB}
\definecolor{green}{HTML}{177C63}
\definecolor{red}{HTML}{BD4935}
\definecolor{orange}{HTML}{AA6C16}
\definecolor{purple}{HTML}{715393}
\tikzset{every node/.style={font=\fontsize{9.3}{12.5}\selectfont,text=ink},
axis/.style={-{Stealth[length=4pt]},line width=.65pt,draw=ink!65},
vec/.style={-{Stealth[length=5pt]},line width=1pt,draw=blue},
short/.style={-{Stealth[length=3pt]},line width=.8pt},
guide/.style={dashed,draw=ink!40,line width=.5pt},
card/.style={rounded corners=2pt,draw=blue!65,fill=blue!4,align=center,text width=6.5cm,inner xsep=9pt,inner ysep=7pt},
title/.style={font=\fontsize{10.5}{14}\selectfont\bfseries,text width=15cm,align=center},
label/.style={font=\fontsize{8.4}{11}\selectfont,align=center},
note/.style={label,text width=6.2cm}}
\begin{document}
\begin{tikzpicture}[x=1cm,y=1cm]
\path[use as bounding box] (-8,-3.5) rectangle(8,3.5);
\node[title] at (0,3.15) {M3-010a\quad 三个高斯面的包围电荷};
\begin{scope}[shift={(-4,0)}]

\foreach \x/\r/\lab in {-1.7/.28/{r<a},0/.66/{a<r<b},1.7/1.0/{r>b}}{
\begin{scope}[shift={(\x,0)},scale=.72]
\filldraw[draw=ink!40,fill=blue!5](0,0)circle(1.1);\filldraw[draw=ink!40,fill=white](0,0)circle(.5);
\fill[red](0,0)circle(2pt);\draw[green,dashed,thick](0,0)circle(\r*1.2);
\end{scope}\node[label] at(\x,-1.3){$\lab$};}
\node[label] at(0,1.6){红点为 $Q$；蓝色为介质球壳；绿圈为高斯面};
\node[label] at(0,-2.2){从左到右：不包围束缚面、只包围内表面、包围两面};
\end{scope}
\node[card,draw=blue!65,fill=blue!4] at (3.9,1.8) {$r<a$：$q_f=Q,\ q_b=0$};
\node[card,draw=green!65,fill=green!4] at (3.9,0) {$a<r<b$：$q_f=Q,\ q_b=-Q\alpha$};
\node[card,draw=orange!65,fill=orange!4] at (3.9,-1.8) {$r>b$：$q_f=Q,\ q_b=-Q\alpha+Q\alpha=0$};
\end{tikzpicture}\end{document}